% Part: sets-functions-relations
% Chapter: infinite
% Section: dedekinds-proof
% Hindi translation (hi-Deva-IN) of Open Logic commit 9620cc73f9c8e0ad003c514a5d3748f29611c4c0.
%
\documentclass[../../../include/open-logic-section]{subfiles}

\begin{document}

\olfileid[hi]{sfr}{infinite}{dedekindsproof}

\olsection[डेडेकिंड का ``प्रमाण'']{अनंत समुच्चय के अस्तित्व के लिए
डेडेकिंड का ``प्रमाण''}

इस अध्याय में हमने डेडेकिंड बीजगणितों के माध्यम से प्राकृत संख्याओं का
समुच्चय-सैद्धांतिक विवेचन प्रस्तुत किया है। \olref[arith][ref]{sec} में हमने
(अन्य बातों के साथ) विश्लेषण के अंकगणितीकरण के दार्शनिक महत्त्व पर विचार किया
था। अब हमें यहाँ प्राप्त उपलब्धि के महत्त्व पर विचार करना चाहिए।

पूरे \olref[sfr][arith][]{chap} में हमने प्राकृत संख्याओं को उपलब्ध मानकर उनका
उपयोग पूर्णांकों, परिमेय संख्याओं और वास्तविक संख्याओं के स्पष्ट निर्माण में
किया। इस अध्याय में हमने प्राकृत संख्याओं का कोई स्पष्ट निर्माण नहीं दिया। हमने
केवल यह दिखाया है कि \emph{डेडेकिंड-अनंत कोई भी समुच्चय दिए जाने पर} हम एक ऐसा
समुच्चय परिभाषित कर सकते हैं जो ठीक उसी तरह व्यवहार करेगा जैसा हम~$\Nat$ से
चाहते हैं।

इसलिए स्पष्टतः हम किसी तत्त्वमीमांसात्मक प्रश्न का उत्तर देने का दावा नहीं कर
सकते, जैसे कि \emph{प्राकृत संख्याएँ कौन-सी वस्तुएँ हैं}। पर यह अच्छी बात है।
आखिरकार, \olref[sfr][arith][ref]{sec} में हमने इस बात पर बल दिया था कि
डेडेकिंड विच्छेदों के समुच्चय के रूप में~$\Real$ की परिभाषा को सुविधाजनक
निर्धारण के बजाय \emph{खोज} समझना गलत होगा। निर्णायक बात यह है कि डेडेकिंड
विच्छेद वास्तविक संख्याओं के प्रमुख गणितीय गुणों को साकार करते हैं। यहाँ भी
ठीक ऐसा ही है: \emph{कोई भी} डेडेकिंड बीजगणित प्राकृत संख्याओं के प्रमुख
गणितीय गुणों को साकार करता है। (वास्तव में डेडेकिंड ने सभी डेडेकिंड बीजगणितों
को \emph{समाकृतिक} सिद्ध करके इस बात को पूरी शक्ति से स्थापित किया
(\citeyear[प्रमेय 132--3]{Dedekind1888})। इसलिए इसमें आश्चर्य नहीं कि आज के कई
``संरचनावादी'' डेडेकिंड को अपना अग्रदूत मानते हैं।)

इसके अतिरिक्त, हमने दिखाया है कि प्राकृत संख्याओं के सिद्धांत को सरल
समुच्चय-सिद्धांत में कैसे अंतःस्थापित किया जा सकता है। वह समुच्चय-सिद्धांत अब
भी कुछ अनौपचारिक है, पर वह (प्रकटतः) प्राकृत संख्याओं को पहले से उपलब्ध नहीं
मानता। अतः शायद हम डेडेकिंड की अपनी महत्त्वाकांक्षी परियोजना को साकार करने की
दिशा में बढ़ रहे हैं। उन्होंने इसे इस प्रकार समझाया था:
\begin{quote}
	विज्ञान में प्रमाणित की जा सकने वाली किसी भी बात को प्रमाण के बिना नहीं
	मानना चाहिए। यद्यपि यह माँग उचित प्रतीत होती है, फिर भी मैं नहीं मान सकता
	कि सबसे सरल विज्ञान की नींव रखने की नवीनतम विधियाँ भी इसे पूरा करती हैं;
	अर्थात् तर्कशास्त्र का वह भाग जो संख्या-सिद्धांत से संबंध रखता है। जब मैं
	अंकगणित (बीजगणित, विश्लेषण) को केवल तर्कशास्त्र का एक भाग कहता हूँ, तो मेरा
	आशय यह है कि संख्या की संकल्पना को मैं स्थान और समय की धारणाओं या
	अंतःप्रज्ञाओं से पूर्णतः स्वतंत्र मानता हूँ---बल्कि उसे चिंतन के शुद्ध
	नियमों का प्रत्यक्ष परिणाम मानता हूँ।
	\citep[प्रस्तावना]{Dedekind1888}
\end{quote}
डेडेकिंड का साहसिक विचार यह है। हमने अभी देखा कि केवल सरल समुच्चय-सिद्धांत के
सहारे प्राकृत संख्याएँ कैसे बनाई जा सकती हैं। \olref[sfr][arith][]{chap} में
हमने देखा कि प्राकृत संख्याएँ और कुछ समुच्चय-सिद्धांत दिए जाने पर वास्तविक
संख्याएँ कैसे बनाई जा सकती हैं। अतः संभव है कि ``अंकगणित (बीजगणित,
विश्लेषण)'', ``तर्कशास्त्र का मात्र एक भाग'' निकले (``तर्कशास्त्र'' शब्द के
डेडेकिंड के विस्तृत अर्थ में)।

विचार तो यही है। पर एक क्षण रुकिए। हमारा डेडेकिंड बीजगणित (प्राकृत संख्याओं का
हमारा स्थानापन्न) तभी बनता है जब कोई डेडेकिंड-अनंत समुच्चय अस्तित्व में हो।
(\olref[sfr][infinite][dedekind]{thm:DedekindInfiniteAlgebra} को फिर से देखें।)
यदि किसी डेडेकिंड-अनंत समुच्चय का अस्तित्व ``तर्कशास्त्र'' या ``चिंतन के शुद्ध
नियमों'' से स्थापित नहीं किया जा सकता, तो परियोजना यहीं रुक जाती है।

तो क्या किसी डेडेकिंड-अनंत समुच्चय का अस्तित्व ``चिंतन के शुद्ध नियमों'' से
स्थापित \emph{किया जा सकता है}? डेडेकिंड ने यह प्रयास किया था:
\begin{quote}
  मेरे विचारों का अपना जगत, अर्थात् उन सभी वस्तुओं की समष्टि~$S$ जो मेरे
  चिंतन का विषय हो सकती हैं, अनंत है। क्योंकि यदि $s$, $S$ के किसी अवयव को
  सूचित करता है, तो यह विचार~$s'$ कि~$s$ मेरे चिंतन का विषय हो सकता है, स्वयं
  $S$ का अवयव है। यदि हम इसे प्रतिबिंब $\phi(s)$, अर्थात् अवयव~$s$ का
  प्रतिबिंब मानें, तो
  \dots~$S$ [डेडेकिंड-]अनंत है; यही सिद्ध करना था।
	\citep[\S66]{Dedekind1888}
\end{quote}
मुख्यतः अत्यंत कठोर गणितीय प्रमाणों से बनी पुस्तक के बीच ऐसी बात मिलना सचमुच
आश्चर्यजनक है। यहाँ दो टिप्पणियाँ करना उचित होगा।

पहली: इस ``प्रमाण'' में लगभग कुछ भी ऐसा नहीं है जिसे आज हम ``गणितीय'' स्वरूप
का मानेंगे। इसमें मनोवैज्ञानिक वस्तुओं (विचारों) की चर्चा होती है, और वह भी
केवल \emph{संभव} विचारों की।

दूसरी: कम-से-कम डेडेकिंड बीजगणितों को हमने जिस रूप में प्रस्तुत किया है, उसमें
इस ``प्रमाण'' की एक सीधी तकनीकी कमी है। यदि डेडेकिंड का तर्क सफल है, तो वह
केवल इतना स्थापित करता है कि अनंत संख्या में वस्तुएँ हैं (विशेषतः, अनंत संख्या
में विचार)। किंतु डेडेकिंड को यह मानने का कारण भी देना होगा कि~$S$ एक अकेला
\emph{समुच्चय} है जिसके !!{element} अनंत संख्या में हैं, न कि~$S$ का आशय
बहुवचन में केवल \emph{कुछ वस्तुओं} से है।

डेडेकिंड ने यहाँ कोई अंतराल नहीं देखा; इससे संकेत मिल सकता है कि उनका
``समष्टि'' शब्द का प्रयोग \emph{हमारे} ``समुच्चय'' शब्द के प्रयोग से ठीक-ठीक
मेल नहीं खाता।\footnote{दरअसल, ऐसा सोचने के हमारे पास दूसरे कारण भी हैं;
\citet[पृ.~23]{Potter2004} देखें।} लेकिन इसमें बहुत आश्चर्य नहीं होगा। पिछले
दो अध्यायों में हमने जो परियोजना अपनाई है---प्राकृत संख्याओं का एक
``निर्माण'', और उनसे पूर्णांकों, वास्तविक संख्याओं तथा परिमेय संख्याओं का
``निर्माण''---वह पूरी तरह सरल ढंग से की गई है। जब जिस समुच्चय की जरूरत पड़ी,
हमने उसे मान लिया; हमने इस बारे में बहुत कम सामान्य सिद्धांत दिए कि ठीक कौन-से
समुच्चय अस्तित्व में हैं और कब। पर हम जानते हैं कि हमें \emph{कुछ} सामान्य
सिद्धांत चाहिए, अन्यथा हम रसेल के विरोधाभास में पड़ जाएँगे।

अब अपने भोलेपन से आगे बढ़ने का समय आ गया है।

\end{document}
